% !TeX root = ../navier-stokes-manuscript.tex
% ============================================================
% 2.2 Finite-Time Response Escalation
% ============================================================

\subsection{Finite-Time Response Escalation}\label{sub:FiniteTimeResponseEscalation}

The Zeno schedule leaves open the only question that matters dynamically: can one solution perform all of those changes before the available time runs out?
A scale can shrink on paper without any fluid moving through the intervening states.
The Navier--Stokes field has no such freedom.
Every finer state must be reached from the state already present, and a finite endpoint shortens the time in which that response can occur.

Along one flow line, define \(A_r=D_t^ru(X,t)\); then \(dA_r/dt=A_{r+1}\).
Its fundamental-theorem identity gives the conditional mechanism needed here: if \(A_r\) escapes by \(T_*\), then \(A_{r+1}\) cannot remain integrable along that path.
This does not make any rung monotone or nonzero; it states what finite-time escape would require inside the same fluid.

The same dependence appears before choosing a particle, but its functional setting must stay fixed throughout the implication.
Fix one Banach space \(B\), retain the Eulerian rows
\[
  V_n(t):=\partial_t^n u(\cdot,t),
\]
and suppose that \(V_n:[0,T_*)\to B\) is locally absolutely continuous, with Bochner derivative \(V_{n+1}\).
For \(0\le s<t<T<T_*\), the Banach-valued fundamental theorem of calculus gives
\[
  V_n(t)-V_n(s)
  =
  \int_s^t V_{n+1}(\tau)\,d\tau
\]
and
\[
  \lVert V_n(t)-V_n(s)\rVert_B
  \le
  \int_s^t \lVert V_{n+1}(\tau)\rVert_B\,d\tau.
\]
If
\[
  \sup_{0<t<T_*}\lVert V_n(t)\rVert_B
  =
  \infty,
\]
then the finiteness of \(\lVert V_n(0)\rVert_B\) forces
\[
  \int_0^{T_*}\lVert V_{n+1}(\tau)\rVert_B\,d\tau
  =
  \infty.
\]
Because \(T_*<\infty\), this in turn forces
\[
  \sup_{0<t<T_*}\lVert V_{n+1}(t)\rVert_B
  =
  \infty.
\]
The implication can be iterated upward only while the same fixed space \(B\), absolute continuity, and Bochner-derivative identity hold at every successive row.
That is the precise Eulerian form of finite-time response escalation.

The material and Eulerian ladders are not identical: already at first order, \eqref{eq:MaterialAccelerationBody} adds the transport derivative, and each further material derivative mixes temporal and spatial derivatives with nonlinear products.
The particle picture supplies the escalation, but the equation must keep the temporal, spatial, and pressure pieces coupled.
A single time index can no longer hold the response, so the next object must retain the full mixed jet.
